%=============================================================================
%  FINAL GROUP PROJECT REPORT
%  ML Model Deployment using Hugging Face and GitHub Actions
%  Software Development for AI Models
%
%  COMPILE WITH: pdflatex (run twice so the Table of Contents fills in)
%
%  FOLDER LAYOUT EXPECTED IN OVERLEAF:
%      report.tex
%      screenshots/      <- the 45 files fig01 ... fig45
%      figures/          <- the 16 files fig_eda_* and fig_model_*
%=============================================================================
\documentclass[11pt,a4paper]{report}

%----------------------------- PAGE GEOMETRY ---------------------------------
\usepackage[a4paper,top=2.6cm,bottom=2.6cm,left=2.6cm,right=2.6cm]{geometry}

%----------------------------- ENCODING & FONTS ------------------------------
\usepackage[T1]{fontenc}
\usepackage[utf8]{inputenc}
\usepackage{lmodern}
\usepackage[english]{babel}
\usepackage{microtype}

%----------------------------- GRAPHICS & TABLES -----------------------------
\usepackage{graphicx}
\usepackage{float}
\usepackage{booktabs}
\usepackage{longtable}
\usepackage{array}
\usepackage{multirow}
\usepackage{caption}
\usepackage{xcolor}
\usepackage{enumitem}
\usepackage{amsmath}
\usepackage{amssymb}

%----------------------------- LINKS -----------------------------------------
\usepackage[hidelinks,breaklinks=true]{hyperref}
\usepackage{url}

%----------------------------- CODE LISTINGS ---------------------------------
\usepackage{listings}

%----------------------------- HEADERS AND FOOTERS ---------------------------
\usepackage{fancyhdr}

%----------------------------- THE PAGE BORDER -------------------------------
% A single thin black rectangle is drawn on EVERY page, including the title
% page and the table of contents. eso-pic places the drawing behind the text
% on every shipped-out page, and TikZ draws the rectangle relative to the
% physical paper edges so it never shifts with the text block.
\usepackage{tikz}
\usetikzlibrary{calc}
\usepackage{eso-pic}

\AddToShipoutPictureBG{%
  \begin{tikzpicture}[remember picture,overlay]
    \draw[line width=0.5pt,black]
      ($(current page.north west)+(1.30cm,-1.30cm)$)
      rectangle
      ($(current page.south east)+(-1.30cm,1.30cm)$);
  \end{tikzpicture}%
}

%----------------------------- STYLE SETTINGS --------------------------------
\definecolor{accentblue}{RGB}{31,78,121}
\definecolor{lightgrey}{RGB}{242,242,242}
\definecolor{codegreen}{RGB}{0,110,60}

\captionsetup{font=small,labelfont=bf,width=0.92\textwidth}
\setlength{\parskip}{0.55em}
\setlength{\parindent}{0pt}
\renewcommand{\arraystretch}{1.25}

% Long code identifiers such as estimated_delivery_days cannot be hyphenated,
% which can push a line past the margin. This gives TeX extra slack to stretch
% inter-word spacing on the affected line rather than overflowing.
\setlength{\emergencystretch}{3em}

\pagestyle{fancy}
\fancyhf{}
\fancyhead[L]{\small\itshape Olist Delivery Risk Radar}
\fancyhead[R]{\small\itshape Software Development for AI Models}
\fancyfoot[C]{\small\thepage}
\renewcommand{\headrulewidth}{0.4pt}
\renewcommand{\footrulewidth}{0pt}

\fancypagestyle{plain}{%
  \fancyhf{}%
  \fancyfoot[C]{\small\thepage}%
  \renewcommand{\headrulewidth}{0pt}%
}

\lstset{
  basicstyle=\ttfamily\scriptsize,
  backgroundcolor=\color{lightgrey},
  frame=single,
  rulecolor=\color{gray},
  breaklines=true,
  breakatwhitespace=true,
  showstringspaces=false,
  keywordstyle=\color{accentblue}\bfseries,
  commentstyle=\color{codegreen}\itshape,
  columns=fullflexible,
  xleftmargin=6pt,
  framexleftmargin=6pt,
  aboveskip=10pt,
  belowskip=10pt
}

% Where to look for images. Overleaf will search both folders in order.
\graphicspath{{screenshots/}{figures/}{./}}

% A shorthand so every figure is laid out identically.
\newcommand{\projfig}[3]{%
  \begin{figure}[H]
    \centering
    \includegraphics[width=#2\textwidth]{#1}
    \caption{#3}
    \label{fig:#1}
  \end{figure}%
}

\begin{document}

%=============================================================================
%  TITLE PAGE
%=============================================================================
\begin{titlepage}
\thispagestyle{empty}
\centering
% The document-wide \parskip would add vertical space after every block here
% and push the closing lines onto a second page. Suppress it locally so the
% whole title page fits on one sheet.
\setlength{\parskip}{0pt}

\vspace*{0.2cm}

{\large\scshape Masters in Artificial Intelligence with Business\par}
\vspace{0.20cm}
{\normalsize Software Development for AI Models\par}

\vspace{0.75cm}
\rule{\textwidth}{1.1pt}
\vspace{0.45cm}

{\huge\bfseries Olist Delivery Risk Radar\par}
\vspace{0.40cm}
{\Large End-to-End Development and Automated\\[0.15cm]
Deployment of a Machine Learning Application\\[0.15cm]
on Hugging Face\par}

\vspace{0.42cm}
\rule{\textwidth}{1.1pt}

\vspace{0.60cm}
{\normalsize\itshape Predicting late deliveries on Brazil's largest e-commerce\\[0.10cm]
marketplace using machine learning, deployed automatically\\[0.10cm]
through GitHub Actions\par}

\vspace{0.85cm}

{\large\bfseries Submitted by\par}
\vspace{0.35cm}
{\normalsize
\begin{tabular}{c}
Kartik Joshi \\[0.16cm]
Prem Kukreja \\[0.16cm]
Samuel Alex \\[0.16cm]
Gagandeep Singh \\[0.16cm]
Aditya Chitale \\[0.16cm]
Karan Baid \\
\end{tabular}\par}

\vspace{0.80cm}

{\large\bfseries Under the guidance of\par}
\vspace{0.28cm}
{\normalsize Prof.\ JP Agarwal\par}

\vspace{0.85cm}

{\normalsize\itshape Submitted for the partial fulfillment of the requirements\\[0.14cm]
for the award of the degree of\par}
\vspace{0.25cm}
{\normalsize\bfseries Master of Science in Artificial Intelligence with Business\par}

\vspace{0.75cm}
{\normalsize Final Group Project\par}
\vspace{0.16cm}
{\normalsize July 2026\par}

\vfill
\end{titlepage}

%=============================================================================
%  TABLE OF CONTENTS
%=============================================================================
\pagenumbering{roman}
\setcounter{page}{1}

\tableofcontents
\clearpage

\listoffigures
\clearpage

\listoftables
\clearpage

\pagenumbering{arabic}
\setcounter{page}{1}

%=============================================================================
\chapter{Problem Statement and Objective}
%=============================================================================

\section{The Business Problem}

Olist is a Brazilian e-commerce marketplace that connects thousands of small,
independent sellers to customers across all twenty-seven Brazilian states. The
seller owns the stock and dispatches the parcel, but the customer buys through
the Olist storefront. This creates a structural weakness that sits at the heart
of this project: \textbf{when a seller delivers late, the customer does not
blame the seller. They blame Olist.}

Our analysis of 96{,}470 completed deliveries quantifies precisely how expensive
that misattributed blame is. Orders that arrived on time received an average
customer review of \textbf{4.29 stars out of 5}. Orders that arrived late
averaged \textbf{2.57 stars}. Most strikingly, the proportion of orders receiving
a damning one-star or two-star review rises from \textbf{9.2\%} when delivery is
punctual to \textbf{54.1\%} when it is late. A late delivery therefore makes a
severely negative review approximately \textbf{5.9 times more likely}.

On a marketplace, poor reviews are not merely embarrassing. They suppress future
search ranking, reduce conversion, and drive customers to competitors. The
damage compounds long after the individual parcel has been delivered.

\section{Objective}

The practical objective of this project is to build a system that identifies
at-risk orders \textbf{at the moment they are placed}, while there is still time
to act.

This timing constraint is the defining engineering requirement of the entire
project. A model that predicts lateness after the parcel has already shipped is
academically interesting and commercially worthless. A model that predicts it at
checkout allows the marketplace to intervene: it can quote a more realistic
delivery date, upgrade the shipping method, prompt the seller to dispatch
immediately, or send a proactive courtesy message that manages the customer's
expectations before those expectations are broken.

In effect, the system predicts a customer-dissatisfaction event at the single
moment when that event is still preventable.

\section{Justification for the Dataset Choice}

The Brazilian E-Commerce Public Dataset by Olist was selected over several
alternatives for five specific reasons, each mapping directly onto a requirement
of this project.

\begin{enumerate}[leftmargin=*]

\item \textbf{It is genuine commercial data, not a teaching set.} The data was
released by Olist itself and covers real transactions between September 2016 and
October 2018. It has been anonymised so thoroughly that company names appearing
in customer review text were replaced with the names of Great Houses from
\textit{Game of Thrones}. This provenance means the data carries the
irregularities of a real operational system rather than the artificial
cleanliness of a synthetic benchmark.

\item \textbf{It contains authentic data-quality problems.} The dataset presents
missing values across thirteen columns, 261{,}831 duplicate rows in the
geolocation table, forty-two GPS coordinates that fall outside Brazil's national
borders, and 775 orders with no line items attached. These are exactly the
issues the project brief requires students to identify and resolve, and they
occur naturally rather than being introduced artificially.

\item \textbf{It presents a genuine data-leakage trap.} The raw data contains
columns such as the actual delivery date and the customer's review score. A
model given these columns would achieve near-perfect scores and be entirely
useless in production, because neither value exists at the moment a prediction
is required. Recognising and removing them is a substantive analytical
contribution rather than a box-ticking exercise.

\item \textbf{It supports meaningful feature engineering.} The dataset includes a
geolocation table mapping Brazilian postcodes to latitude and longitude. Joining
this table to both the seller and the customer allowed us to compute the
physical distance a parcel must travel, a variable that does not exist anywhere
in the raw files and which proved to be one of the model's most informative
inputs.

\item \textbf{It has an unambiguous business case with quantifiable value.} As
established in Section 1.1, the link between late delivery and customer
dissatisfaction is measurable and severe. This allows the value of the model to
be argued in commercial terms rather than purely statistical ones.

\end{enumerate}

\section{How the Solution Adds Value}

The deployed system converts a passive historical record into an active
operational instrument. Three concrete benefits follow.

\textbf{Targeted intervention rather than blanket policy.} Without a model, a
marketplace can only apply uniform rules, such as padding every delivery estimate
by three days. That degrades the customer experience for the 91.9\% of orders
that would have arrived on time anyway. A per-order risk score allows resources
to be concentrated on the minority of orders that actually need attention.

\textbf{Quantified prioritisation.} The model returns a calibrated probability
rather than a binary label. An order scored at 36\% is not merely ``risky'', it
is 4.4 times riskier than the marketplace average, and that ratio can be used to
rank a day's orders and allocate a limited operations budget rationally.

\textbf{A measurable reduction in expected cost.} As documented in Section 5.5,
the decision threshold was selected to minimise expected business cost. Against
the current baseline of taking no action at all, the deployed configuration
reduces expected cost by \textbf{28.9\%} on the held-out test set.

%=============================================================================
\chapter{Dataset Description and Source}
%=============================================================================

\section{Source and Licence}

\begin{table}[H]
\centering
\caption{Dataset provenance}
\begin{tabular}{@{}ll@{}}
\toprule
\textbf{Attribute} & \textbf{Value} \\
\midrule
Title & Brazilian E-Commerce Public Dataset by Olist \\
Publisher & Olist Store (official public release) \\
Platform & Kaggle \\
URL & \url{https://www.kaggle.com/datasets/olistbr/brazilian-ecommerce} \\
Licence & CC BY-NC-SA 4.0 \\
Period covered & September 2016 to October 2018 \\
Total orders & 99{,}441 \\
Tables & 9 relational CSV files \\
Raw size & Approximately 120 MB \\
\bottomrule
\end{tabular}
\end{table}

\projfig{fig01_dataset_source}{0.95}{The dataset source page on Kaggle, showing
the official Olist release, the usability rating and the CC BY-NC-SA 4.0
licence under which the data is distributed.}

\section{Structure of the Nine Tables}

The data arrives as nine separate files in a normalised relational structure.
Every table was audited before any transformation was applied.

\begin{table}[H]
\centering
\caption{Inventory of the nine raw Olist tables}
\begin{tabular}{@{}lrrl@{}}
\toprule
\textbf{Table} & \textbf{Rows} & \textbf{Columns} & \textbf{Contents} \\
\midrule
geolocation & 1{,}000{,}163 & 5 & Postcode to latitude/longitude mapping \\
order\_items & 112{,}650 & 7 & One row per physical item purchased \\
payments & 103{,}886 & 5 & Payment method and instalments \\
customers & 99{,}441 & 5 & Customer location \\
orders & 99{,}441 & 8 & Order status and five timestamps \\
reviews & 99{,}224 & 7 & Customer review scores and text \\
products & 32{,}951 & 9 & Category, weight and dimensions \\
sellers & 3{,}095 & 4 & Seller location \\
category\_translation & 71 & 2 & Portuguese to English category names \\
\midrule
\textbf{Total} & \textbf{1{,}550{,}922} & & \\
\bottomrule
\end{tabular}
\end{table}

\projfig{fig03_audit_inventory}{0.98}{Automated audit of the nine raw tables,
reporting row counts, column counts and memory footprint, followed by the
structure and status breakdown of the central \texttt{orders} table.}

\section{Definition of the Target Variable}

An order is defined as \textbf{late} when the date the customer actually
received it falls after the delivery date they were promised at checkout:

\begin{equation}
\texttt{is\_late} =
\begin{cases}
1 & \text{if } \texttt{order\_delivered\_customer\_date} > \texttt{order\_estimated\_delivery\_date} \\
0 & \text{otherwise}
\end{cases}
\end{equation}

Only orders with a status of \texttt{delivered} and with both dates present can
be labelled, since an order still in transit has no known outcome and a
cancelled order was never delivered at all. This reduces the dataset as follows:

\begin{table}[H]
\centering
\caption{Construction of the modelling population}
\begin{tabular}{@{}lrr@{}}
\toprule
\textbf{Filtering step} & \textbf{Orders} & \textbf{Retained} \\
\midrule
All orders in the dataset & 99{,}441 & 100.0\% \\
Status equal to \texttt{delivered} & 96{,}478 & 97.0\% \\
Both delivery dates present & 96{,}470 & 97.0\% \\
Joined successfully to order items & 96{,}470 & 97.0\% \\
\bottomrule
\end{tabular}
\end{table}

Of these 96{,}470 orders, \textbf{7{,}826 (8.11\%)} arrived late. When an order
is late, it is late by a median of 5.8 days, though the distribution has a long
tail extending to 189 days.

\projfig{fig05_audit_target}{0.92}{Verification that all four foreign-key
relationships join at 100\%, that all primary keys are unique, and the resulting
class balance of the target variable.}

%=============================================================================
\chapter{Exploratory Data Analysis}
%=============================================================================

\section{Data Quality Audit}

The audit was performed as a read-only pass over the raw files, deliberately
separated from any cleaning, so that decisions could be made with full knowledge
of the data rather than reactively.

\begin{table}[H]
\centering
\caption{Data-quality issues identified in the raw data}
\begin{tabular}{@{}lll@{}}
\toprule
\textbf{Issue} & \textbf{Scale} & \textbf{Resolution} \\
\midrule
Missing review comment titles & 87{,}656 (88.3\%) & Column not used \\
Missing review comment text & 58{,}247 (58.7\%) & Column not used \\
Missing delivery dates & 2{,}965 (3.0\%) & Rows excluded, undeliverable \\
Missing product attributes & 610 (1.9\%) & Imputed in pipeline \\
Missing carrier dates & 1{,}783 (1.8\%) & Column excluded as leaky \\
Duplicate geolocation rows & 261{,}831 (26.2\%) & Aggregated by median \\
GPS points outside Brazil & 42 & Removed before aggregation \\
Orders with no line items & 775 (0.78\%) & Excluded by inner join \\
Order-value outliers & 7{,}658 (7.9\%) & Retained, genuine purchases \\
\bottomrule
\end{tabular}
\end{table}

\projfig{fig04_audit_missing}{0.9}{Missing-value profile across all nine tables
and the duplicate-row count per table. The 26.2\% duplication in the
geolocation table is expected: it stores many GPS readings per postcode.}

\projfig{fig_eda_11_missing_values}{0.82}{Missing values remaining in the
assembled analytical table. All are handled by imputation inside the modelling
pipeline rather than by deleting rows.}

\section{Key Observations}

Ten observations emerged from the analysis. Each is reproduced below with the
supporting evidence.

\subsection{The problem is imbalanced, so accuracy is misleading}

Only 8.11\% of delivered orders arrived late. A model that ignores its inputs
entirely and always answers ``on time'' would achieve \textbf{91.9\% accuracy}
while catching zero late deliveries. Accuracy is therefore an actively
misleading headline metric for this problem, a point returned to in
Section 5.2.

\projfig{fig_eda_01_target_balance}{0.9}{Class distribution of the target
variable. The rarity of the positive class is the single most important
statistical property of this problem.}

\subsection{Distance is strongly predictive}

Orders travelling under 50 km are late 6.5\% of the time. Orders travelling
beyond 1{,}500 km are late \textbf{13.1\%} of the time, roughly double the rate.
This column did not exist in the raw data and was engineered from the
geolocation table, as described in Section 4.1.

\projfig{fig_eda_02_distance}{0.95}{The engineered distance feature. The left
panel shows the distribution of journey lengths split by outcome; the right
panel shows late rate rising monotonically with distance.}

\subsection{A tight promise is a broken promise}

Orders promised within seven days are late \textbf{31.3\%} of the time, against
only 5.6\% for orders given between 31 and 60 days. The promised delivery window
is legitimate as a model input because it is displayed to the customer on the
checkout page and is therefore genuinely known at prediction time.

\projfig{fig_eda_03_promise_window}{0.95}{Delivery performance against the length
of the promise made at checkout. A generous promise is easy to keep; an
aggressive one is not.}

\subsection{Performance is highly unstable over time}

The monthly late rate swings from 1.1\% in October 2016 to \textbf{21.4\%} in
March 2018, a nineteen-fold difference around an overall average of 8.1\%. Two
distinct disruptions are visible: a sharp Black Friday spike in November 2017
reaching 14.3\%, roughly triple the preceding months, and a sustained
deterioration through February and March 2018 whose cause is not recorded in
the dataset. Because performance depends so heavily on when an order is placed,
purchase month and week of year were provided to the model.

\projfig{fig_eda_04_monthly_trend}{0.95}{Order volume and late-delivery rate by
month. Order volume grows steadily while reliability fluctuates violently.}

\subsection{Delivery reliability varies sharply by region}

Alagoas records the highest late rate at \textbf{23.9\%}, while Rond\^onia is the
most reliable at 2.9\%. States geographically distant from the S\~ao Paulo seller
hub perform consistently worse, which is consistent with the distance finding in
Section 3.2.2.

\projfig{fig_eda_05_state_analysis}{0.95}{Late rate by destination state, and the
relationship between average journey distance and reliability. Bubble size
represents order volume.}

\subsection{Product category matters}

The \texttt{audio} category is late 13.1\% of the time against only 5.4\% for
\texttt{luggage\_accessories}. Bulky and fragile categories are over-represented
at the risky end of the distribution.

\projfig{fig_eda_06_category_analysis}{0.82}{Riskiest and safest product
categories among those with at least 300 orders.}

\subsection{Money and size columns are heavily right-skewed}

Order value has a median of R\$~86 but a maximum of R\$~13{,}440. In total 7{,}658
orders (7.9\%) fall outside the conventional 1.5$\times$IQR range. These were
deliberately retained rather than removed: they represent genuine high-value
purchases, and the tree-based models selected for this problem are insensitive
to monotonic skew.

\projfig{fig_eda_07_distributions}{0.98}{Distribution of the principal numeric
features, with the top 1\% trimmed for readability.}

\projfig{fig_eda_08_outliers}{0.95}{Boxplots confirming the long right tail
present in every money and size column.}

\subsection{No single feature explains lateness on its own}

The strongest linear relationships with the target are all weak:

\begin{table}[H]
\centering
\caption{Strongest linear correlations with the target}
\begin{tabular}{@{}lr@{}}
\toprule
\textbf{Feature} & \textbf{Correlation} \\
\midrule
\texttt{distance\_km} & $+0.069$ \\
\texttt{estimated\_delivery\_days} & $-0.060$ \\
\texttt{same\_state} & $-0.056$ \\
\bottomrule
\end{tabular}
\end{table}

This is an important diagnostic
finding rather than a disappointing one. Weak individual correlations predict
that a linear model will struggle and that tree-based ensembles, which capture
interactions between features, should perform substantially better. Section 5.3
confirms this prediction directly.

\projfig{fig_eda_09_correlation}{0.78}{Correlation matrix of numeric features
against the target.}

\subsection{The commercial justification}

Orders delivered on time average 4.29 stars; late orders average 2.57. The share
of one and two-star reviews rises from 9.2\% to \textbf{54.1\%}, an increase of
5.9 times. This single relationship justifies the entire project.

\projfig{fig_eda_10_review_impact}{0.95}{The commercial consequence of late
delivery. The left panel shows the full distribution of review scores by
outcome; the right panel isolates the share of severely negative reviews.}

\subsection{Residual missing values are handled by imputation}

After joining, fifteen columns still contain gaps, the largest being
\texttt{avg\_photos\_qty} at 1.38\%. All are handled by imputation inside the
modelling pipeline, never by deleting rows, so that no observation is discarded
for a missing value in a secondary attribute.

\projfig{fig08_eda_complete}{0.98}{Completion of the exploratory analysis, with
the commercial finding and confirmation that eleven figures were written to
disk.}

%=============================================================================
\chapter{Data Preprocessing}
%=============================================================================

\section{Assembling One Table from Nine}

A machine-learning model requires a single flat table in which each row is one
observation. The nine relational tables were therefore joined into one
analytical table of \textbf{96{,}470 orders by 53 columns}.

The join was not a simple merge. The \texttt{order\_items} table contains one row
per physical item, so an order containing three products appears three times. To
produce one row per order, item-level data was aggregated: item counts, distinct
seller counts, total price, total freight, total weight and maximum item value.
For attributes that cannot be summed, such as the product category, the value
was taken from the \textbf{most expensive item} in the order, on the reasoning
that this item dominates packaging, handling and routing.

\projfig{fig07_build_dataset}{0.98}{Construction of the analytical table:
geolocation cleaning, the four-way join, distance computation and the resulting
96{,}470 by 53 table.}

\section{Feature Engineering}

Fifteen features were derived that do not exist in the raw data.

\subsection{The distance feature}

The most substantial contribution is \texttt{distance\_km}, the physical distance
between seller and buyer. Constructing it required three steps.

First, the geolocation table's 1{,}000{,}163 GPS readings were filtered to
Brazil's true bounding box, removing forty-two coordinates that placed Brazilian
postcodes in other continents. Second, the surviving readings were collapsed to
one representative point per postcode using the \textbf{median} rather than the
mean, so that a small number of erroneous readings could not drag a postcode's
position. This reduced one million rows to 19{,}010 unique postcodes. Third, the
great-circle distance between the seller's postcode and the customer's postcode
was computed using the haversine formula:

\begin{equation}
d = 2r \arcsin\left(\sqrt{\sin^{2}\left(\frac{\phi_2-\phi_1}{2}\right) +
\cos\phi_1 \cos\phi_2 \sin^{2}\left(\frac{\lambda_2-\lambda_1}{2}\right)}\right)
\end{equation}

where $\phi$ denotes latitude, $\lambda$ denotes longitude and $r$ is the Earth's
radius of 6{,}371 km. The haversine formula is required rather than simple
subtraction because the Earth is a sphere: one degree of longitude represents a
very different ground distance near the poles than at the equator.

Distance was successfully computed for \textbf{99.51\%} of orders, with a median
journey of 434 km and a maximum of 3{,}398 km.

\subsection{Remaining engineered features}

\begin{table}[H]
\centering
\caption{Features engineered from the raw data}
\begin{tabular}{@{}lll@{}}
\toprule
\textbf{Group} & \textbf{Feature} & \textbf{Rationale} \\
\midrule
Geography & \texttt{distance\_km} & Physical journey length \\
 & \texttt{same\_state} & Intra-state deliveries avoid transfers \\
 & \texttt{is\_sao\_paulo\_seller} & Proximity to the logistics hub \\
\midrule
Time & \texttt{purchase\_month} & Seasonal courier congestion \\
 & \texttt{purchase\_week\_of\_year} & Captures Black Friday precisely \\
 & \texttt{purchase\_day\_of\_week} & Weekday dispatch patterns \\
 & \texttt{purchase\_hour} & Cut-off times for same-day dispatch \\
 & \texttt{purchase\_is\_weekend} & Reduced weekend operations \\
\midrule
Promise & \texttt{estimated\_delivery\_days} & Length of the commitment made \\
\midrule
Money & \texttt{freight\_ratio} & Shipping cost relative to value \\
 & \texttt{price\_per\_item} & Average item value \\
 & \texttt{freight\_per\_kg} & Shipping service level proxy \\
\midrule
Logistics & \texttt{multi\_seller\_order} & Multiple dispatch points \\
 & \texttt{weight\_per\_item\_g} & Handling difficulty \\
\bottomrule
\end{tabular}
\end{table}

\section{Preventing Data Leakage}

\subsection{What data leakage is}

Data leakage occurs when a model is given information that would not be
available at the moment the prediction is required. The analogy is a student who
sees the examination answers in advance: the resulting score is excellent and
entirely meaningless, because the conditions under which it was obtained will
never recur.

\subsection{The governing rule}

A single rule was applied throughout: \textbf{the model may only use facts the
marketplace knows at the instant the customer clicks Buy.} Any column recorded
afterwards is prohibited without exception.

\begin{table}[H]
\centering
\caption{Columns excluded to prevent data leakage}
\begin{tabular}{@{}p{6.2cm}p{7.2cm}@{}}
\toprule
\textbf{Column excluded} & \textbf{Reason} \\
\midrule
\texttt{order\_delivered\_customer\_date} & The target is constructed from it \\
\texttt{order\_delivered\_carrier\_date} & Recorded days after checkout \\
\texttt{order\_approved\_at} & Recorded once payment clears \\
\texttt{order\_estimated\_delivery\_date} & Retained only as a derived duration \\
\texttt{order\_status} & Known only once the order concludes \\
\texttt{actual\_delivery\_days} & Derived from the actual delivery date \\
\texttt{days\_late} & Literally the answer \\
\texttt{review\_score} & Written by the customer after delivery \\
\texttt{purchase\_year} & Would not generalise to future orders \\
identifier columns & Random codes carry no signal \\
city columns & Over 4{,}000 values would explode the encoding \\
postcode columns & Latitude and longitude capture location better \\
\bottomrule
\end{tabular}
\end{table}

\projfig{fig06_audit_leakage}{0.95}{The data-leakage policy, generated
automatically by the audit script so that the exclusions are documented in code
rather than asserted in prose.}

\subsection{A subtle case worth stating explicitly}

The feature \texttt{estimated\_delivery\_days}, the number of days between
purchase and the promised delivery date, might superficially appear to be
leakage because the target is defined relative to that same promised date. It is
not. The promised date is calculated by the marketplace and displayed to the
customer on the checkout page \textit{before} payment is taken. It is therefore
genuinely available at prediction time. The question the model answers is
whether the promise that was made will be kept, and the ambitiousness of that
promise is legitimately part of the answer.

\projfig{fig09_feature_policy}{0.95}{The complete feature policy: excluded
columns with reasons, the thirty numeric and four categorical features retained,
and the explicit note regarding the promised delivery window.}

\section{The Preprocessing Pipeline}

All preprocessing is implemented as a scikit-learn \texttt{Pipeline} rather than
as free-standing transformations. Numeric and categorical columns follow
separate paths that execute in parallel through a \texttt{ColumnTransformer}.

\begin{table}[H]
\centering
\caption{Preprocessing applied by column type}
\begin{tabular}{@{}lll@{}}
\toprule
\textbf{Step} & \textbf{Numeric (30 columns)} & \textbf{Categorical (4 columns)} \\
\midrule
Missing values & Median imputation & Constant, ``missing'' \\
Transformation & Standard scaling & One-hot encoding \\
Unseen values & Not applicable & Grouped as infrequent \\
Rare categories & Not applicable & Merged below 30 occurrences \\
\bottomrule
\end{tabular}
\end{table}

Two choices merit explanation.

\textbf{Median imputation} was preferred to the mean because the money columns
are heavily right-skewed, and a mean would be dragged upward by a small number
of very large orders.

\textbf{One-hot encoding was configured to tolerate unseen categories.} The
encoder is created with the following option:

\begin{lstlisting}[language=Python]
OneHotEncoder(handle_unknown="infrequent_if_exist", min_frequency=30)
\end{lstlisting}

This ensures the deployed application cannot crash when a user selects a product
category the model never saw during training; such values are routed to a
``rare'' bucket instead of raising an exception. This behaviour is verified
explicitly by the smoke test described in Section 7.5.

\section{Train and Test Split}

The data was divided into 80\% training and 20\% testing using a
\textbf{stratified} split, which preserves the 8.11\% positive rate in both
halves so that the test set is neither easier nor harder than the training set.

\begin{table}[H]
\centering
\caption{Train and test split}
\begin{tabular}{@{}lrr@{}}
\toprule
\textbf{Partition} & \textbf{Orders} & \textbf{Late rate} \\
\midrule
Training & 77{,}176 & 8.11\% \\
Test (held out) & 19{,}294 & 8.11\% \\
\midrule
\textbf{Total} & \textbf{96{,}470} & \textbf{8.11\%} \\
\bottomrule
\end{tabular}
\end{table}

Critically, \textbf{every preprocessing step is fitted on the training partition
only}. The medians used for imputation, the means and standard deviations used
for scaling, and the category vocabulary used for encoding are all learned from
training data and then merely applied to the test data. Because these steps live
inside the pipeline object, this separation is enforced structurally rather than
depending on the analyst remembering to maintain it. During cross-validation the
entire pipeline is refitted within each fold, so no fold is ever influenced by
another fold's statistics.

\projfig{fig10_train_test_split}{0.95}{The stratified split and the explicit
statement of how leakage is prevented at the preprocessing stage.}

%=============================================================================
\chapter{Model Training and Evaluation}
%=============================================================================

\section{Models Selected}

Three models were trained. The selection is deliberately one model per
\textbf{algorithm family}, so that the comparison tests genuinely different
approaches rather than padding the results table with near-identical variants.

\begin{table}[H]
\centering
\caption{Models compared and the rationale for each}
\begin{tabular}{@{}lp{4.1cm}p{5.6cm}@{}}
\toprule
\textbf{Model} & \textbf{Family} & \textbf{Why it earns its place} \\
\midrule
Logistic Regression & Linear & The honest baseline. Without it, no claim about
the value of the ensembles can be substantiated \\
Random Forest & Bagging & Independent trees vote; handles non-linear
interactions and provides feature importances \\
Gradient Boosting & Boosting & Sequential error correction; typically the
strongest family on tabular data \\
\bottomrule
\end{tabular}
\end{table}

Three further models were considered and \textbf{deliberately rejected}, and the
reasoning is recorded here because a defensible exclusion is as informative as
an inclusion. A support vector machine with a radial basis kernel would require
hours to train on 77{,}176 rows and would not be expected to outperform gradient
boosting. K-nearest neighbours degrades badly in the sixty-plus dimensions
produced by one-hot encoding, a well-understood consequence of the curse of
dimensionality. A standalone decision tree is strictly dominated by the random
forest, which is itself an ensemble of decision trees. Adding any of these would
have lengthened the results table without adding information.

All three models were configured with \texttt{class\_weight="balanced"} so that
the rare late-delivery class is weighted in proportion to its scarcity.
Without this, all three models converge on the degenerate solution of always
predicting ``on time''.

\section{Evaluation Metrics and Why They Are Appropriate}

Because only 8.11\% of orders are late, metric selection is not a formality. The
following metrics were computed, and each answers a different question.

\begin{table}[H]
\centering
\caption{Evaluation metrics and their interpretation}
\begin{tabular}{@{}lp{9.3cm}@{}}
\toprule
\textbf{Metric} & \textbf{Question it answers} \\
\midrule
Accuracy & What proportion of all predictions were correct? Reported for
completeness only; a constant ``on time'' model scores 91.9\% \\
Precision & Of the orders we flagged, how many really were late? Low precision
means wasted intervention effort \\
Recall & Of the orders that really were late, how many did we catch? Low recall
means dissatisfied customers we never saw coming \\
F1 & The harmonic mean of precision and recall \\
ROC-AUC & How well does the model \textit{rank} orders from safe to risky across
all possible cut-offs? Used as the primary model-selection metric \\
PR-AUC & As above but focused solely on the rare positive class, making it the
more honest summary under class imbalance \\
Brier score & How far are the printed probabilities from observed reality?
Used to evaluate calibration in Section 5.4 \\
\bottomrule
\end{tabular}
\end{table}

\textbf{ROC-AUC was selected as the primary metric} because the operational task
is prioritisation: given a day's orders and a limited operations budget, which
should be examined first? That is a ranking problem, and ROC-AUC measures ranking
quality independently of any particular threshold.

\textbf{Recall is treated as the primary business metric.} The asymmetry is
established directly by the exploratory analysis. A missed late delivery
attracts a one or two-star review 54.1\% of the time. A false alarm costs one
courtesy email. Failing to detect a genuine problem is materially more expensive
than investigating one that turns out to be fine.

\section{Performance Comparison}

All three models were assessed on the same held-out 19{,}294 orders that none of
them saw during training, and additionally under five-fold stratified
cross-validation on the training partition.

\begin{table}[H]
\centering
\caption{Model performance on the held-out test set}
\label{tab:comparison}
\begin{tabular}{@{}lrrrrrrr@{}}
\toprule
\textbf{Model} & \textbf{Acc.} & \textbf{Prec.} & \textbf{Recall} & \textbf{F1}
& \textbf{ROC-AUC} & \textbf{PR-AUC} & \textbf{CV AUC} \\
\midrule
\textbf{Gradient Boosting} & 0.777 & 0.219 & 0.686 & 0.332 & \textbf{0.810} &
\textbf{0.336} & 0.806 \\
Random Forest & 0.845 & 0.273 & 0.551 & 0.365 & 0.799 & 0.330 & 0.794 \\
Logistic Regression & 0.657 & 0.144 & 0.655 & 0.236 & 0.710 & 0.183 & 0.712 \\
\midrule
\textit{Always ``on time''} & \textit{0.919} & --- & \textit{0.000} & --- &
\textit{0.500} & \textit{0.081} & --- \\
\bottomrule
\end{tabular}
\end{table}

Table \ref{tab:comparison} contains the most important single line in this
report, and it is the last one. \textbf{A model that ignores every input and
always answers ``on time'' achieves 91.9\% accuracy, higher than any of our three
trained models, while catching precisely zero late deliveries.} This is why
accuracy was not used to select the model, and why any report quoting accuracy
alone on an imbalanced problem should be treated with suspicion.

The cross-validation standard deviation was below 0.006 for all three models,
indicating that the performance differences are stable rather than artefacts of
a fortunate split.

\projfig{fig11_model_comparison}{0.95}{Model comparison as generated by the
training script, including the degenerate baseline for reference.}

\projfig{fig_model_01_roc_pr_curves}{0.95}{ROC and precision-recall curves for
all three models. The precision-recall curve is the more informative of the two
under class imbalance, since it ignores the abundant negative class.}

The exploratory analysis predicted this outcome. Section 3.2.8 observed that no
individual feature correlates strongly with the target, from which it follows
that a linear model would struggle and that ensembles capturing feature
interactions would perform substantially better. The measured gap between
logistic regression (0.710) and the two ensembles (0.799 and 0.810) confirms
that prediction quantitatively.

\section{Probability Calibration}

\subsection{The problem discovered}

The winning model ranked orders well but printed probabilities that were simply
wrong. Because \texttt{class\_weight="balanced"} instructs the model to treat the
rare class as though it were common, every predicted probability was inflated.

The average predicted probability across the test set was \textbf{36.0\%} against
a true late rate of \textbf{8.11\%}, an overstatement of roughly
\textbf{4.4 times}. Inspected by band, orders the model scored between 30\% and
40\% were in reality late only 3.4\% of the time.

This mattered because the project brief requires the application to display
prediction probability. Displaying a number that is wrong by a factor of four
would have been misleading to the user.

\subsection{The correction applied}

\textbf{Isotonic regression} was fitted using five-fold cross-validation via
\texttt{CalibratedClassifierCV}. Isotonic regression learns a monotonic mapping
from the model's raw scores onto the frequencies actually observed. Because the
mapping is monotonic, the ranking of orders is preserved exactly, which is why
ROC-AUC is unaffected.

\begin{table}[H]
\centering
\caption{Effect of isotonic calibration}
\begin{tabular}{@{}lrrr@{}}
\toprule
\textbf{Measure} & \textbf{Raw model} & \textbf{Calibrated} & \textbf{Truth} \\
\midrule
Mean predicted probability & 0.3600 & \textbf{0.0813} & 0.0811 \\
Brier score (lower is better) & 0.1595 & \textbf{0.0633} & --- \\
ROC-AUC & 0.8099 & 0.8101 & --- \\
PR-AUC & 0.3355 & 0.3354 & --- \\
\bottomrule
\end{tabular}
\end{table}

The mean predicted probability moved to within 0.0002 of the true rate and the
Brier score fell by 60\%, while ROC-AUC and PR-AUC were unchanged to three
decimal places. Discrimination was preserved and honesty was gained.

\projfig{fig12_calibration}{0.92}{Detection and correction of the probability
inflation.}

\projfig{fig_model_05_calibration}{0.95}{Reliability curve before and after
calibration. The calibrated line follows the diagonal closely, meaning a
displayed probability of 20\% corresponds to approximately one order in five
genuinely arriving late.}

\section{Selecting the Decision Threshold}

A classifier outputs a probability, not a decision. Converting one into the
other requires a threshold, and that choice is a business decision rather than a
statistical one.

Two conventional options were rejected. The default of 0.5 is arbitrary and
particularly inappropriate when one class is rare. Maximising the F1 score is
also arbitrary, because F1 implicitly assumes that a missed late delivery and a
false alarm carry equal cost, which directly contradicts the evidence in
Section 3.2.9.

Instead the threshold was chosen to minimise expected business cost:

\begin{equation}
\text{Cost}(t) = R \times \text{FN}(t) + 1 \times \text{FP}(t)
\end{equation}

where $t$ is the threshold, FN and FP are false negatives and false positives at
that threshold, and $R$ expresses how many false alarms one missed late delivery
is worth. Rather than invent a currency figure, a sensitivity analysis was
performed across plausible values of $R$.

\begin{table}[H]
\centering
\caption{Sensitivity of the optimal threshold to the assumed cost ratio}
\begin{tabular}{@{}lrrr@{}}
\toprule
\textbf{Cost ratio $R$} & \textbf{Optimal threshold} & \textbf{Recall} &
\textbf{Precision} \\
\midrule
3 : 1 & 0.23 & 0.380 & 0.386 \\
\textbf{5 : 1 (shipped)} & \textbf{0.15} & \textbf{0.532} & \textbf{0.304} \\
10 : 1 & 0.12 & 0.613 & 0.267 \\
\bottomrule
\end{tabular}
\end{table}

\begin{table}[H]
\centering
\caption{Comparison of candidate thresholds at $R=5$}
\begin{tabular}{@{}lrrrrr@{}}
\toprule
\textbf{Strategy} & \textbf{Threshold} & \textbf{Prec.} & \textbf{Recall} &
\textbf{F1} & \textbf{Cost} \\
\midrule
Flag nothing (current practice) & --- & --- & 0.000 & --- & 7{,}825 \\
Default 0.50 & 0.50 & 0.639 & 0.069 & 0.125 & 7{,}346 \\
Maximum F1 & 0.18 & 0.332 & 0.479 & 0.392 & 5{,}590 \\
\textbf{Minimum cost (shipped)} & \textbf{0.15} & \textbf{0.304} &
\textbf{0.532} & \textbf{0.387} & \textbf{5{,}563} \\
\bottomrule
\end{tabular}
\end{table}

The shipped threshold of \textbf{0.155} reduces expected cost by \textbf{28.9\%}
relative to taking no action. Note that maximising F1 would have selected a
higher cut-off and caught meaningfully fewer late deliveries, which is the
correct choice only under the assumption that a false alarm is as damaging as a
dissatisfied customer.

\projfig{fig13_threshold_choice}{0.95}{Threshold selection driven by expected
business cost, with sensitivity analysis across three cost ratios.}

\projfig{fig_model_04_threshold_choice}{0.95}{The cost curve is notably flat near
its minimum, indicating the choice is robust to moderate errors in the assumed
cost ratio.}

\projfig{fig_model_02_confusion_matrix}{0.95}{Confusion matrices at the default
and tuned thresholds.}

\section{Justification for the Selected Model}

\textbf{Gradient Boosting was selected}, on the following grounds.

\textbf{Highest ranking quality.} At 0.810 ROC-AUC it leads Random Forest (0.799)
and decisively outperforms Logistic Regression (0.710). Since the operational
task is prioritisation, ranking quality is the correct primary criterion.

\textbf{Highest recall at equivalent operating points.} At the default threshold
it achieves 0.686 recall against Random Forest's 0.551, and recall is the metric
the business asymmetry favours.

\textbf{Stability.} Cross-validation standard deviation of 0.0035 is the lowest
of the three models.

\textbf{Deployment characteristics.} The serialised pipeline occupies 1.14 MB,
comfortably below Hugging Face's inline file limit, and returns a prediction in
well under a second on free CPU hardware.

\textbf{Native handling of missing values.} \texttt{HistGradientBoostingClassifier}
handles gaps internally, providing a second line of defence behind the
pipeline's imputation.

\projfig{fig_model_03_feature_importance}{0.9}{Permutation importance for the
selected model, measured as the drop in ROC-AUC when each column is randomly
shuffled. The engineered \texttt{distance\_km} feature appears among the top six
inputs despite not existing in the raw data.}

\section{Honest Limitations}

\textbf{Precision is 0.304.} Roughly seven in every ten flagged orders would in
fact have arrived on time. This is acceptable only because the intervention is
inexpensive; a courtesy email costs far less than the one-star review that 54\%
of genuinely late deliveries attract. Should intervention costs rise, the
threshold should be raised with them, and Section 5.5 provides the analysis
needed to do so.

\textbf{The model reflects 2016 to 2018 courier conditions.} Logistics networks
change, and the model would require periodic retraining in genuine operational
use.

\textbf{The February to March 2018 disruption is unexplained.} The model learns
the pattern without understanding its cause, and would not anticipate a similar
future event arising from different circumstances.

\textbf{Seller-level history was excluded.} A seller's past punctuality would
almost certainly improve predictions, but computing it safely requires
time-aware aggregation to avoid reintroducing leakage. It is recorded here as
future work rather than attempted unsafely.

\textbf{ROC-AUC of 0.810 indicates good ranking, not certainty.} The system is
decision support, not an oracle.

%=============================================================================
\chapter{Model Saving}
%=============================================================================

The final model was serialised with \texttt{joblib}, which is preferred to plain
\texttt{pickle} for scikit-learn objects because it stores large NumPy arrays
efficiently.

\begin{lstlisting}[language=Python,caption={Saving the complete pipeline}]
import joblib
joblib.dump(best_model, "model.joblib", compress=3)
\end{lstlisting}

The critical design decision is \textbf{what} was saved. Rather than storing the
classifier alone and reimplementing preprocessing inside the application, the
entire pipeline was saved as one object:

\begin{center}
\texttt{imputer $\rightarrow$ scaler $\rightarrow$ one-hot encoder $\rightarrow$
gradient boosting $\rightarrow$ isotonic calibrator}
\end{center}

This has three consequences. The application requires no preprocessing code of
its own, so preprocessing cannot drift out of step with the model. Deployment
involves one file rather than four that must be kept mutually consistent. And
the exact transformations applied at training time are guaranteed to be applied
at prediction time, because they are literally the same fitted objects.

Alongside the model, \texttt{model\_metadata.json} records the decision
threshold, the feature list, all performance metrics and the exact library
versions used for training, so that the deployed application never hard-codes a
number that could silently diverge from the model.

\begin{table}[H]
\centering
\caption{Saved model artefacts}
\begin{tabular}{@{}llr@{}}
\toprule
\textbf{File} & \textbf{Contents} & \textbf{Size} \\
\midrule
\texttt{model.joblib} & Complete fitted pipeline & 1.14 MB \\
\texttt{model\_metadata.json} & Threshold, features, metrics, versions & 4 KB \\
\bottomrule
\end{tabular}
\end{table}

\projfig{fig14_model_saved}{0.95}{The per-class classification report at the
shipped threshold, followed by model serialisation.}

\projfig{fig16_model_file}{0.85}{The saved artefacts on disk.}

%=============================================================================
\chapter{Application Development}
%=============================================================================

\section{Architecture}

The application is a single-page web application backed by a Python API, served
from one process.

\begin{center}
\begin{tabular}{@{}c@{}}
\fbox{\begin{minipage}{0.86\textwidth}
\centering
\vspace{2mm}
\textbf{Browser} \\
React 19 + TypeScript, deck.gl WebGL map, GLSL shader background \\
$\downarrow$ \texttt{fetch("/api/...")} $\uparrow$ \\
\textbf{Server} \\
FastAPI + Uvicorn (\texttt{app.py}) \\
$\downarrow$ \\
\textbf{Model} \\
\texttt{model.joblib}: full pipeline with isotonic calibration
\vspace{2mm}
\end{minipage}}
\end{tabular}
\end{center}

The model executes server-side because browsers cannot run Python, and because
this keeps the model file itself off the user's machine. The frontend is
compiled to static files at build time and served by the same FastAPI process,
so the entire application occupies one container listening on one port.

\begin{table}[H]
\centering
\caption{API endpoints}
\begin{tabular}{@{}lll@{}}
\toprule
\textbf{Endpoint} & \textbf{Method} & \textbf{Purpose} \\
\midrule
\texttt{/api/config} & GET & Dropdown options, route distances, model metadata \\
\texttt{/api/dashboard} & GET & Filtered aggregates over all 96{,}470 orders \\
\texttt{/api/predict} & POST & Score one order \\
\texttt{/api/health} & GET & Liveness check used by the container \\
\bottomrule
\end{tabular}
\end{table}

\section{The Dashboard}

The dashboard presents analytics across the full historical dataset, with four
filters that update every visualisation simultaneously. Aggregation is performed
server-side, so the browser receives a few hundred numbers rather than a
three-megabyte data file.

\projfig{fig18_app_dashboard_top}{0.98}{The dashboard: filters, five key
performance indicators, the monthly trend and the map of Brazil.}

\projfig{fig19_app_dashboard_map}{0.98}{The WebGL choropleth rendered with
deck.gl. Hovering any state reveals its late rate, order volume, typical journey
distance and revenue. Colour encodes late rate; circle size encodes volume.}

\projfig{fig20_app_dashboard_charts}{0.98}{The three risk drivers and the
commercial impact panel, reproducing the exploratory findings interactively.}

\projfig{fig21_app_dashboard_filtered}{0.98}{Filtering demonstrated. Restricting
to the three remote northern states raises the observed late rate from 8.11\% to
14.65\% and the typical journey from 434 km to 2{,}290 km. Unselected states are
dimmed on the map.}

\section{The Prediction Interface}

The prediction form collects twelve inputs and derives the thirty-four features
the model requires. This asymmetry is deliberate: a user should not be asked to
supply a freight-to-weight ratio when it can be computed from values they have
already entered.

Four presets are provided that land in each of the four risk bands, allowing the
full behaviour of the model to be demonstrated in four clicks.

\projfig{fig22_predict_low}{0.98}{A low-risk order: a local delivery with a
generous twenty-day promise scores 1.4\%, which is 0.17 times the marketplace
average. All three drivers are favourable.}

\projfig{fig23_predict_high}{0.98}{A high-risk order: the same local route, but
promised in six days, scores 22.3\%. The tight promise alone moves the order
across the action threshold.}

\projfig{fig24_predict_veryhigh}{0.98}{A very high-risk order: S\~ao Paulo to
Amazonas, 2{,}693 km, promised in twelve days, during Black Friday week. The
score of 36.0\% is 4.44 times the marketplace average, and the interface
explains each contributing factor.}

\section{Model Transparency}

A dedicated tab exposes how the model was built, including the full comparison
table, feature importances, the leakage policy and the honest limitations. This
was included so that the application does not present itself as an
unquestionable authority.

\projfig{fig25_app_model_insights}{0.98}{The model insights tab: headline
metrics, the three-model comparison and permutation importances.}

\projfig{fig26_app_leakage_limitations}{0.98}{The leakage policy, the calibration
explanation and the honest statement of the model's limitations, all surfaced to
the end user.}

\projfig{fig27_app_about}{0.98}{Project background, the six-step method, dataset
provenance, technology stack and team.}

\section{Batch Prediction}

\subsection{Why a second prediction mode was needed}

The single-order form suits a call-centre agent checking one customer. It is
useless to an operations team facing a morning's five thousand orders. A model
that can only be consulted one record at a time is a lookup tool; a model that
can score a whole day's trading before any of it ships is an operational
instrument.

A second mode was therefore added: upload a CSV, score every row in one pass,
download the results.

\subsection{Design}

The CSV columns are deliberately identical to the twelve fields the web form
collects, so anyone who can use the form already understands the file format. A
correctly formatted template of twelve example orders is downloadable from
within the application. Its first four rows are the four presets from the
single-order form; the remaining eight were chosen to span the input space ---
all four payment types, same-state and cross-country routes, single and
multi-seller orders, parcels from 300~g to 12~kg, and promises from eight to
eighteen days --- so that a user opening the template sees the full range of
inputs the endpoint accepts. Scored together they fall across all four risk
bands, five low through four very high.

\begin{table}[H]
\centering
\caption{The batch prediction endpoints}
\begin{tabular}{@{}lll@{}}
\toprule
\textbf{Endpoint} & \textbf{Method} & \textbf{Purpose} \\
\midrule
\texttt{/api/predict/batch} & POST & Score an uploaded CSV \\
\texttt{/api/predict/batch/template} & GET & Download a formatted example \\
\bottomrule
\end{tabular}
\end{table}

Three design decisions are worth recording.

\textbf{Single and batch prediction share one function.} Both paths call the
same \texttt{build\_features} routine that converts twelve user inputs into the
thirty-four model columns. Had they been implemented separately, the two would
eventually have drifted and begun returning different answers for the same
order. Sharing the function makes that impossible by construction, and it was
verified directly: the first four template rows score 1.4\%, 11.7\%, 36.0\% and
22.3\%, exactly matching the four presets on the single-order form. Keeping
those four rows at the head of the template in form order is deliberate, so the
check can be repeated by eye at any time.

\textbf{Rows are validated individually.} A single malformed line reports its
own row number and is skipped; the rest of the file still scores. The
alternative, rejecting the entire upload because one date was mistyped, would
be infuriating on a five-thousand-row file.

\textbf{The whole batch is scored in one vectorised call} rather than a loop of
single predictions. On two thousand orders this is the difference between
0.4~seconds and several minutes.

\subsection{Output}

Each row gains six columns: the computed journey distance, the calibrated
probability as a decimal and as a percentage, the risk tier, a binary flag, the
multiple of average risk, and the recommended action. Any \texttt{order\_id},
\texttt{customer\_id}, \texttt{reference} or \texttt{notes} column present in
the upload is carried through untouched so results can be joined back to the
operations team's own records.

\projfig{fig49_batch_upload}{0.98}{Batch prediction. Five hundred orders scored
in one pass, with the risk breakdown across the four bands, revenue behind the
flagged orders quantified, and two deliberately malformed rows reported by line
number while the remaining rows scored normally.}

\subsection{Programmatic access}

Because the feature is an HTTP endpoint rather than a browser-only control, the
same functionality is available to a scheduled job:

\begin{lstlisting}[language=Python,caption={Scoring a file from Python}]
import requests

with open("orders.csv", "rb") as fh:
    response = requests.post(
        "https://samuelalex37-olist-delivery-risk.hf.space/api/predict/batch",
        files={"file": fh},
    )

result = response.json()
print(result["summary"])
\end{lstlisting}

Every endpoint is documented interactively at \texttt{/docs}, generated
automatically from the server code so the documentation cannot fall out of date
with the implementation.

A note on framework choice. The reference material for this feature proposed
Gradio, on the reasonable grounds that a Gradio application is automatically
exposed as an API. That benefit does not apply here: FastAPI already publishes
an OpenAPI specification and interactive documentation for every route.
Introducing Gradio would have added a second user-interface framework alongside
React to obtain something the existing stack already provided.

\section{Local Testing}

Before any deployment, the application was tested locally and the saved model was
verified independently by an automated smoke test covering four scenarios.

\begin{table}[H]
\centering
\caption{Smoke test coverage}
\begin{tabular}{@{}lll@{}}
\toprule
\textbf{Test} & \textbf{Verifies} & \textbf{Result} \\
\midrule
Load & \texttt{model.joblib} deserialises & Pass \\
Predict & Accepts a hand-written order & Pass \\
Rank & Risky order scores above safe order & Pass, 36.9\% vs 0.3\% \\
Robustness & Unseen category does not crash & Pass \\
\bottomrule
\end{tabular}
\end{table}

The fourth test is the most important for deployment safety. It submits a
product category and a state code that do not exist in the training data,
confirming that the encoder routes them to its ``rare'' bucket rather than
raising an exception that would return a server error to a live user.

\projfig{fig17_app_terminal}{0.95}{The application running locally on port
7860.}

\projfig{fig15_smoke_test}{0.95}{All four smoke tests passing against the saved
model.}

%=============================================================================
\chapter{Deployment on Hugging Face Spaces}
%=============================================================================

\section{Choice of Space Type}

Hugging Face Spaces offers Gradio, Streamlit, Static and Docker runtimes. The
\textbf{Docker} SDK was selected because the application is a FastAPI server
paired with a compiled React bundle, a combination none of the pre-built
runtimes support.

Configuration is declared in the YAML frontmatter of \texttt{README.md}, which
Hugging Face reads on every push:

\begin{lstlisting}[caption={Space configuration in README.md}]
---
title: Olist Delivery Risk Radar
emoji:
colorFrom: blue
colorTo: indigo
sdk: docker
app_port: 7860
license: cc-by-nc-sa-4.0
short_description: Predicting late delivery risk for Brazilian e-commerce
---
\end{lstlisting}

\projfig{fig34_hf_space_create_form}{0.9}{Creating the Space with the Docker SDK
on free CPU hardware and public visibility.}

\section{The Container Build}

A two-stage Dockerfile is used. The first stage uses a Node image purely to
compile the React application. The second stage begins from a clean Python image
and copies only the compiled output across, so that Node and its 146 npm
packages never reach the runtime image.

\begin{table}[H]
\centering
\caption{Files deployed to the Space}
\begin{tabular}{@{}llr@{}}
\toprule
\textbf{File} & \textbf{Purpose} & \textbf{Size} \\
\midrule
\texttt{app.py} & FastAPI backend and static server & 20.6 KB \\
\texttt{model.joblib} & Trained pipeline & 1.14 MB \\
\texttt{model\_metadata.json} & Threshold, features, metrics & 4 KB \\
\texttt{app\_config.json} & Route distances, dropdown options & 9.8 KB \\
\texttt{dashboard\_data.parquet} & 96{,}470 orders, 16 columns & 2.93 MB \\
\texttt{requirements.txt} & Pinned Python dependencies & 1 KB \\
\texttt{README.md} & Space configuration and documentation & 7.8 KB \\
\texttt{Dockerfile} & Build instructions & 3.7 KB \\
\midrule
\textbf{Total} & & \textbf{4.11 MB} \\
\bottomrule
\end{tabular}
\end{table}

Dependency versions are pinned exactly. A model serialised by scikit-learn 1.9.0
may refuse to deserialise under a different version, so pinning guarantees the
Space runs precisely the libraries that trained the model.

\projfig{fig36_hf_manual_push}{0.95}{The first push to the Space, uploading all
eighteen Git LFS objects.}

\projfig{fig37_hf_build_logs}{0.98}{The Docker build completing: React compiled
in 5.83 seconds, Python dependencies installed, model and data copied, image
pushed.}

\section{Verification of the Deployed Application}

\begin{table}[H]
\centering
\caption{Verification against the live URL}
\begin{tabular}{@{}lll@{}}
\toprule
\textbf{Check} & \textbf{Expected} & \textbf{Observed} \\
\midrule
\texttt{/api/health} & Model name and row count & Gradient Boosting, 96{,}470 \\
\texttt{/api/dashboard} & Aggregates match analysis & 8.11\% late, 27 states \\
\texttt{/api/predict} & Matches local prediction & 36.0\%, 4.44$\times$, 2{,}692.7 km \\
Frontend & Renders without errors & No console errors \\
\bottomrule
\end{tabular}
\end{table}

The prediction returned by the live Space is \textbf{identical} to the one
produced locally, confirming that the pipeline survived serialisation, Git LFS
transfer, container build and network delivery without alteration.

\projfig{fig39_hf_space_files}{0.98}{The Space file listing, showing all deployed
files and the eleven-commit history mirrored from GitHub.}

\projfig{fig38_hf_space_live}{0.98}{The deployed application running publicly on
Hugging Face Spaces, with the browser address bar visible. The status indicator
beside the Space name reads \emph{Running}.}

%=============================================================================
\chapter{GitHub Repository}
%=============================================================================

\section{Repository Structure}

\begin{lstlisting}[caption={Repository layout}]
olist-delivery-risk/
|
+-- app.py                      FastAPI backend and static file server
+-- model.joblib                Complete pipeline (Git LFS)
+-- model_metadata.json         Metrics, threshold, feature list
+-- app_config.json             State coordinates, route distances
+-- dashboard_data.parquet      96,470 orders, 16 columns (Git LFS)
+-- requirements.txt            Deployment dependencies, pinned
+-- requirements-dev.txt        Adds plotting and the Streamlit prototype
+-- Dockerfile                  Two-stage build
+-- .dockerignore
+-- .gitattributes              LF line endings, LFS tracking
+-- .gitignore
+-- README.md                   Space config and documentation
+-- streamlit_app.py            Earlier prototype, retained as fallback
|
+-- frontend/                   React + TypeScript source
|   +-- src/components/         Dashboard, Predict, ModelInsights,
|   |                           About, BrazilMap, ShaderBackground
|   +-- src/assets/             Simplified Brazil GeoJSON (118 KB)
|   +-- package.json, vite.config.ts, tsconfig.json
|
+-- src/                        Machine-learning pipeline, in run order
|   +-- 01_data_audit.py        Inspect the nine raw tables
|   +-- 02_build_dataset.py     Join tables, engineer features
|   +-- 03_eda.py               Eleven figures, ten observations
|   +-- 04_train_models.py      Train, calibrate, threshold, save
|   +-- 05_smoke_test.py        Verify the saved model
|   +-- 06_build_app_assets.py  Compress data for deployment
|   +-- 07_simplify_geojson.py  Douglas-Peucker map simplification
|   +-- utils.py                Shared paths and helpers
|
+-- reports/figures/            Sixteen analysis figures
+-- data/raw/                   Kaggle CSVs (gitignored, 120 MB)
|
+-- .github/workflows/
    +-- deploy.yml              Automated deployment
\end{lstlisting}

Two exclusions are deliberate. The 120 MB of raw Kaggle data is not committed,
since data belongs in a data store rather than in version control and the README
documents how to obtain it. Compiled frontend output is not committed either,
because it is regenerated by every build and a stale committed bundle could
otherwise ship by accident.

\projfig{fig30_github_repo_root}{0.98}{The public repository.}

\projfig{fig32_github_workflows_folder}{0.9}{The workflow directory.}

\section{Version Control Practice}

The project history comprises eleven commits, each representing a coherent stage
of work with a message explaining the reasoning rather than restating the diff.

\projfig{fig31_github_commit_history}{0.98}{Commit history.}

\projfig{fig29_git_push_output}{0.95}{Initial push to GitHub.}

Binary files are tracked with \textbf{Git LFS}. This became necessary rather than
optional: Hugging Face rejects any push containing binary files inline. Eighteen
files were migrated, rewriting history so that binaries never appear inline in
any commit.

%=============================================================================
\chapter{Automated Deployment with GitHub Actions}
%=============================================================================

\section{How the Automated Deployment Works}

GitHub Actions is an automation service built into GitHub. A workflow file
describes a checklist, and GitHub executes that checklist automatically on a
chosen trigger. The complete deployment sequence is:

\begin{enumerate}[leftmargin=*]
\item A developer pushes a commit to the \texttt{main} branch.
\item GitHub detects the push and starts the workflow.
\item \textbf{Job 1, validate.} A fresh Ubuntu machine installs the pinned
Python dependencies, loads \texttt{model.joblib} and runs the smoke test, starts
the FastAPI server and polls \texttt{/api/health} until it answers, then installs
the npm packages and compiles the React application. Any failure stops the
workflow here.
\item \textbf{Job 2, deploy.} Only if every validation step passed, a second
machine checks out the full history including LFS content, reads the Hugging
Face token from the encrypted secret store, uploads the LFS objects and force
pushes the repository to the Space.
\item Hugging Face detects the new commit, rebuilds the Docker image and
restarts the container.
\item The updated application is live, typically within four minutes of the
original push.
\end{enumerate}

No human touches Hugging Face at any point after the initial Space creation.

\projfig{fig40_workflow_file}{0.98}{The workflow file in the repository.}

\section{Why Validation Precedes Deployment}

The two-job structure is the most consequential design decision in the pipeline.
A single-job workflow that merely pushed files would deploy a broken application
to a public URL, where the failure would be discovered by whoever opened it
next.

By gating deployment behind \texttt{needs: validate}, a defect fails on the
build machine and the live Space is left untouched, still serving the last
known-good version. The practical value of this is demonstrated directly by the
run history in Section 9.4.

\section{Secure Credential Handling}

The project brief requires that the Hugging Face token must not be hard-coded.
Three measures were applied.

\textbf{The token is stored as an encrypted GitHub Actions secret.} Once saved,
it cannot be read back by anyone, including the person who created it. The
interface displays only its name.

\textbf{The token is fine-grained.} Rather than a general write token with
authority over the entire Hugging Face account, a scoped token was issued with
read and write permission on \textbf{one repository only}. Even in the event of
compromise, the blast radius is a single Space that can be rebuilt in minutes.
This is the principle of least privilege.

\textbf{The token never appears in the repository.} The workflow references it
only as \texttt{\$\{\{ secrets.HF\_TOKEN \}\}}, and GitHub automatically masks the
value in all logs. A search of the complete repository history confirms no
token literal has ever been committed.

\projfig{fig35_github_secret_added}{0.92}{The encrypted secret. Only the name and
last-updated time are visible; the value cannot be retrieved.}

\section{Testing Automatic Deployment}

To satisfy Task 4.4, a deliberately visible change was made: a green badge
reading ``Auto-deployed by GitHub Actions'' was added to the application header.
The change was chosen so that its presence on the live site is unambiguous
evidence that the pipeline works.

\projfig{fig41_visible_change_commit}{0.98}{The commit containing the visible
change.}

\projfig{fig42_actions_run_success}{0.98}{The successful workflow execution. Both
jobs completed in one minute and eight seconds, and the custom job summary
reports the destination Space.}

\projfig{fig43_actions_deploy_steps}{0.98}{The deployment job in detail. The
expanded step shows eighteen LFS objects uploaded and the commit range
\texttt{613b08e..a55447c} pushed to the Space. The credential is masked
automatically.}

\projfig{fig44_hf_space_updated}{0.98}{The live Space carrying the new badge. No
person uploaded this change to Hugging Face.}

\begin{table}[H]
\centering
\caption{Timeline of the automated deployment}
\begin{tabular}{@{}ll@{}}
\toprule
\textbf{Time} & \textbf{Event} \\
\midrule
14:41 & Commit \texttt{a55447c} pushed to \texttt{main} \\
14:42 & Workflow run \#6 begins \\
14:43 & Both jobs complete successfully (1m 8s) \\
14:44 & Space serves the newly compiled bundle \\
14:45 & Badge visible on the public URL \\
\bottomrule
\end{tabular}
\end{table}

Objective confirmation is provided by the compiled bundle filename, which changed
from \texttt{index-BmV02qjo.js} to \texttt{index-DaqtONgO.js}. The Space is
serving a file that did not exist before the push and that no person uploaded.

\section{Evidence That the Validation Gate Is Real}

The Actions history contains six runs, and the pattern is more informative than
a single success would be.

\begin{table}[H]
\centering
\caption{Complete workflow run history}
\begin{tabular}{@{}clllc@{}}
\toprule
\textbf{Run} & \textbf{Commit} & \textbf{Validate} & \textbf{Deploy} &
\textbf{Reason} \\
\midrule
1 & \texttt{d57ea57} & Pass & Fail & No secret, Space not yet created \\
2 & \texttt{90668fb} & Pass & Fail & Secret still absent \\
3 & \texttt{ae0b503} & Pass & Fail & Secret still absent \\
4 & \texttt{eafd62d} & Pass & Fail & Secret still absent \\
5 & \texttt{613b08e} & Pass & Fail & Secret still absent \\
\textbf{6} & \texttt{a55447c} & \textbf{Pass} & \textbf{Pass} &
\textbf{Secret configured} \\
\bottomrule
\end{tabular}
\end{table}

Validation passed on every single run, including the five that failed to deploy.
Every failure occurred at the deployment step and for a correctly diagnosed
reason. This demonstrates that the gate is functional rather than decorative:
the pipeline consistently distinguished ``the code is sound'' from ``the
deployment credentials are absent''.

\projfig{fig45_actions_history}{0.98}{The full run history, showing the
progression from failing deployments to a complete success once the secret was
configured.}

%=============================================================================
\chapter{The Automated MLOps Pipeline}
%=============================================================================

\section{Why a Second Workflow Exists}

The workflow described in Chapter 10 deploys the \textit{application}. It does
not retrain the \textit{model}. Those are different activities with different
risk profiles, and conflating them causes two problems.

First, retraining on every commit would produce a slightly different model each
time. Every performance figure quoted in this report would then stop matching
the live application, and a reader checking the deployed system against the
document would find they disagree.

Second, most commits change the frontend. Attaching a fifteen-minute training
run to a two-line CSS change wastes both time and the repository's allocation of
compute minutes.

The project therefore has two workflows with clearly separated
responsibilities.

\begin{table}[H]
\centering
\caption{The two workflows and their responsibilities}
\begin{tabular}{@{}lll@{}}
\toprule
\textbf{Workflow} & \textbf{Trigger} & \textbf{Responsibility} \\
\midrule
\texttt{deploy.yml} & Every push to \texttt{main} & Ships the application \\
\texttt{mlops-pipeline.yml} & Manual or scheduled & Retrains and registers the model \\
\bottomrule
\end{tabular}
\end{table}

\section{The Four Stages}

The pipeline follows the structure taught in the course material, adapted to
this project's data and accounts.

\begin{center}
\begin{tabular}{@{}c@{}}
\fbox{\begin{minipage}{0.9\textwidth}
\centering
\vspace{2mm}
\textbf{Stage 1 · Register dataset} \\
publish the analytical table to a Hugging Face \textit{dataset} repository \\
$\downarrow$ \\
\textbf{Stage 2 · Data preparation} \\
freeze a stratified train/test split and publish it \\
$\downarrow$ \\
\textbf{Stage 3 · Model training} \\
train three models, calibrate the winner, register it as a \textit{model} \\
$\downarrow$ \\
\textbf{Stage 4 · Promote to Space} \\
run pre-flight checks, then publish to the live application
\vspace{2mm}
\end{minipage}}
\end{tabular}
\end{center}

Each stage runs on its own freshly provisioned machine. They share no memory and
no filesystem, so they pass work to one another through the Hugging Face Hub,
which acts as the pipeline's versioned storage layer.

\begin{table}[H]
\centering
\caption{Hugging Face repositories used by the pipeline}
\begin{tabular}{@{}lll@{}}
\toprule
\textbf{Repository} & \textbf{Type} & \textbf{Contents} \\
\midrule
\texttt{olist-delivery-risk-data} & Dataset & Analytical table and frozen splits \\
\texttt{olist-delivery-risk-model} & Model & Trained pipeline and metadata \\
\texttt{olist-delivery-risk} & Space & The running application \\
\bottomrule
\end{tabular}
\end{table}

Separating the three means the model can be retrained and re-registered without
touching the application, and the application can be redeployed without
retraining. Every upload creates a commit, so a model can always be traced back
to the exact data it was trained on.

Both repositories publish a card describing their contents, and both cards are
generated by the pipeline rather than written by hand, so they cannot drift out
of date with the artefacts they describe.

The dataset card also solves a concrete problem. The Hub's viewer concatenates
every Parquet file it finds in a repository, and this repository holds the full
analytical table alongside the train and test partitions derived from it. The
viewer therefore reported 192,940 rows for a dataset containing 96,470 orders,
having counted the same data once whole and once split. Declaring three explicit
configurations in the card resolves it: each view now reports its own correct
count.

\projfig{fig47_hf_dataset_repo}{0.9}{The dataset card. Each configuration
reports its own row count, and the note explains why viewing them together would
double-count.}

\projfig{fig48_hf_model_repo}{0.98}{The model registry. The card is generated
from the metrics of the run that produced the model, and the Hub has
automatically linked the dataset it was trained on and the Space that serves it.}

\section{Freezing the Split}

Stage 2 splits the data once and publishes the result, rather than each training
run making its own split. This matters more than it first appears.

If every run split independently, two models could not be compared fairly: a
difference in score might reflect nothing more than one model receiving an
easier test set. Publishing the split converts model comparison from an anecdote
into a controlled experiment.

The stage also enforces the leakage policy in code. Before splitting anything it
asserts that none of the nine banned columns appears in the feature list, and
terminates the pipeline if one does. This is a hard failure rather than a
warning, because a pipeline that trains a cheating model and reports excellent
scores is worse than one that stops.

\section{Experiment Tracking}

Stage 3 logs every model, parameter and metric to MLflow. Because a GitHub
runner is destroyed when its job finishes, the tracking database would normally
vanish with it. The workflow therefore uploads it as a build artefact, so the
experiment history outlives the machine that produced it and can be downloaded
and browsed afterwards:

\begin{lstlisting}[caption={Inspecting the tracking history from a pipeline run}]
mlflow ui --backend-store-uri sqlite:///mlflow.db
\end{lstlisting}

One implementation detail is worth recording, because it was not the obvious
choice. MLflow 3.14 places the familiar \texttt{./mlruns} filesystem store into
maintenance mode and refuses to open it, so tracking uses a SQLite database
instead. That is still a single self-contained file requiring no server, but on
the supported code path.

An honest limitation: this is per-run history rather than a permanent central
tracking server. A production team would host MLflow separately so that every
run across every branch accumulates in one queryable place.

\section{Gated Promotion}

Stage 4 is the only stage that changes what users see, and it is deliberately
the most cautious. It runs only when explicitly requested, and performs three
checks before doing anything irreversible.

\begin{enumerate}[leftmargin=*]
\item \textbf{The model loads.} A corrupted or mismatched artefact fails here
rather than on the live site.
\item \textbf{It predicts sensibly.} A deliberately risky order must score above
a deliberately safe one. A model that has these the wrong way round is broken
regardless of what its metrics claim.
\item \textbf{It is not materially worse than what is already deployed.} A drop
of more than 0.02 ROC-AUC against the live model stops the promotion.
\end{enumerate}

The third check is the important one. A pipeline that automatically ships
whatever it most recently trained will eventually ship a regression. Requiring
the candidate to justify replacing the incumbent turns automation from a risk
into a safeguard.

\section{Reproducibility}

The pipeline was executed three times: once during the original development,
once locally through the new scripts, and once on a GitHub Actions runner from a
clean checkout. All three produced identical results.

\begin{table}[H]
\centering
\caption{The same model reproduced across three independent executions}
\begin{tabular}{@{}lccc@{}}
\toprule
\textbf{Metric} & \textbf{Original} & \textbf{Pipeline (local)} &
\textbf{Pipeline (CI)} \\
\midrule
Logistic Regression ROC-AUC & 0.7097 & 0.7097 & 0.7097 \\
Random Forest ROC-AUC & 0.7994 & 0.7994 & 0.7994 \\
Gradient Boosting ROC-AUC & 0.8099 & 0.8099 & 0.8099 \\
Mean probability before calibration & 36.00\% & 36.00\% & 36.00\% \\
Mean probability after calibration & 8.13\% & 8.13\% & 8.13\% \\
Brier score after calibration & 0.0633 & 0.0633 & 0.0633 \\
Decision threshold & 0.155 & 0.155 & 0.155 \\
Expected cost reduction & 28.9\% & 28.9\% & 28.9\% \\
\midrule
Operating system & Windows & Windows & Ubuntu \\
\bottomrule
\end{tabular}
\end{table}

This is the strongest evidence in the report that the modelling is sound. The
figures were not produced by one fortunate run on one machine. They are
reproducible from the registered data, on a different operating system, through
a different code path, by a process with no human involvement.

Training took one minute and thirty-eight seconds on the runner against roughly
ten minutes locally, because the hosted machine has more cores available to the
parallel forest and the five-fold calibration.

\projfig{fig46_mlops_pipeline_run}{0.98}{A complete pipeline execution on GitHub
Actions. Stages 1 to 3 completed successfully; stage 4 shows as skipped because
promotion was not requested, which is the intended default.}

%=============================================================================
\chapter{Challenges and Resolutions}
%=============================================================================

Nine substantive problems arose during development. Each is recorded with its
diagnosis and resolution, because the resolutions are more instructive than a
narrative in which nothing went wrong. One is recorded together with the two
incorrect diagnoses that preceded it, since how a fault was misread is part of
what the fault has to teach.

\section{Predicted Probabilities Were Inflated Fourfold}

\textbf{Problem.} Class weighting caused the model to overstate every
probability by roughly 4.4 times. Orders scored between 30\% and 40\% were in
reality late only 3.4\% of the time.

\textbf{Why it mattered.} The brief requires the application to display
prediction probability. Displaying a number wrong by a factor of four would
have misinformed the user.

\textbf{Resolution.} Isotonic calibration, documented in Section 5.4. The Brier
score improved by 60\% while ROC-AUC remained unchanged.

\section{The Threshold Logic Contradicted the Evidence}

\textbf{Problem.} The initial implementation selected the threshold by
maximising F1, while the accompanying analysis argued that recall matters more
than precision. These two positions are mutually inconsistent, because F1 treats
false negatives and false positives as equally costly.

\textbf{Resolution.} Replaced with explicit cost minimisation and a sensitivity
analysis across three cost ratios, documented in Section 5.5.

\section{Hugging Face Rejected Binary Files}

\textbf{Problem.} The first push to the Space was rejected because it contained
binary files inline. Six files exceeded the platform's threshold.

\textbf{Resolution.} Eighteen binary files were migrated to Git LFS, rewriting
history so that no commit contains an inline binary.

\textbf{A consequential secondary defect.} The migration introduced a fault that
would have broken deployment silently. \texttt{actions/checkout} does not fetch
LFS content by default; it downloads small text pointer files instead. The build
machine would therefore have attempted to load a model file whose contents were
a few lines of text, and the Space would have received a placeholder rather than
a working model. The fault was identified before deployment and corrected by
setting \texttt{lfs: true} on both checkout steps and by uploading LFS objects
explicitly, so that any future failure reports plainly rather than manifesting
as a corrupt model.

\section{Four Dependencies That Only Failed on a Clean Machine}

Four separate faults shared one cause: a package was present in the development
environment as an indirect dependency, so the code worked locally while the
declared requirements were incomplete. Each surfaced only on a freshly
provisioned machine.

\textbf{python-multipart was missing.} FastAPI parses uploaded files through
this library and raises at import time if a route declares a file parameter
without it. The application did not merely lose the batch feature; it refused to
start at all. The validation job caught it, the deployment job was skipped, and
the live Space continued serving the previous working version rather than a
container that would crash on boot. This is the two-job structure described in
Section 10.2 paying for itself in a way that could not be arranged deliberately.

\textbf{MLflow could not coexist with pandas 3.} The full MLflow package caps
pandas below version 3, while the application is pinned to 3.0.5. Installing
both is an outright \texttt{ResolutionImpossible}. The cap originates in
MLflow's bundled tracking server, which this pipeline does not run, so the fix
was to depend on \texttt{mlflow-skinny} - the client alone. Logging behaves
identically.

\textbf{MLflow retired the filesystem tracking store.} Version 3.14 places the
familiar \texttt{./mlruns} directory into maintenance mode and refuses to open
it, so the training stage failed on its first tracking call. Tracking was moved
to the supported SQLite backend, which is still a single self-contained file
requiring no server.

\textbf{SQLAlchemy and Alembic were unpinned.} The SQLite backend needs both,
\texttt{mlflow-skinny} does not install them, and they were present locally only
as residue from the full MLflow package that had already been uninstalled. Both
are now declared explicitly.

The common lesson is that a working development machine proves very little about
a requirements file. Only a clean install proves it, which is precisely what a
continuous-integration job provides.

\section{A Credential Was Captured in a Screenshot}

\textbf{Problem.} During documentation, a screenshot inadvertently captured the
Hugging Face access token in plain text at the moment of creation.

\textbf{Resolution.} The exposure was identified during a review pass. A search
of the complete repository history confirmed the token had never been committed
and the screenshots directory is excluded from version control, so the exposure
was limited to a single local file. The token was revoked, a replacement issued,
and the screenshot removed from the documentation set. The incident is recorded
here because credential handling is an explicit requirement of the brief, and
because the correct response to an exposed secret is to revoke it rather than to
hope it goes unnoticed.

\section{A Rate Limit That Presented as an Authentication Failure}

\textbf{Problem.} A deployment failed at the step that uploads binaries to the
Space. The job reported \texttt{LFS: Rate limit exceeded} against Hugging Face's
LFS batch endpoint. The upload command carried \texttt{-{}-all}, which offers
every binary object in the repository on every run, and several pushes had
landed within the same afternoon.

\textbf{Why it mattered.} The validation job had passed, so the repository
recorded a successful build while the Space silently remained a commit behind.
A deployment that is advertised as automatic but quietly stops deploying is
worse than one that never claimed to be automatic, because nothing prompts
anyone to look.

\textbf{The diagnosis was wrong twice before it was right.} The failure was
first attributed to a revoked access token, then to a token whose repository
scope had been narrowed. Both hypotheses were plausible: a token had genuinely
been revoked earlier in the project, and a push from a developer machine
succeeded while the same push from the runner failed. Neither was correct. The
local push had used a different, broader credential held in the operating
system's credential store, so the comparison that appeared to isolate the CI
token was measuring two different secrets. The job log named the true cause in
one line and was decisive from the outset. Reasoning from indirect evidence
while the direct evidence sat unread cost more time than the defect itself.

\textbf{Resolution.} The upload is retried twice with a growing pause. A rate
limit is a request to wait rather than grounds to abandon a deployment. If all
three attempts fail the step still exits non-zero, and it does so before the
reference is pushed, so a deployment whose binaries did not upload is never
presented as successful. The \texttt{-{}-all} flag was dropped in favour of
naming the branch, which offers what the deployment actually requires; on a
single-branch repository this is a modest reduction, and the retry is what makes
the step survive. Lock verification was disabled because Hugging Face does not
implement the LFS locking API and the resulting warning appeared immediately
above the lines that matter when the step fails. The retry logic was exercised
against a substituted \texttt{git} covering success at the first attempt,
success after two failures, and exhaustion of all three. The fix was confirmed
by a subsequent commit that deployed automatically end to end.

%=============================================================================
\chapter{Conclusion}
%=============================================================================

\section{Summary}

This project delivered a complete machine-learning system, from nine raw
relational tables to a publicly accessible application that redeploys itself
automatically.

\begin{table}[H]
\centering
\caption{Summary of outcomes}
\begin{tabular}{@{}ll@{}}
\toprule
\textbf{Dimension} & \textbf{Outcome} \\
\midrule
Dataset & 96{,}470 orders assembled from nine tables \\
Engineered features & 15, including haversine seller-buyer distance \\
Models compared & 3, one per algorithm family \\
Selected model & Gradient Boosting, ROC-AUC 0.810 \\
Calibration & Isotonic; Brier improved 60\% \\
Threshold & 0.155, minimising expected cost \\
Expected cost reduction & 28.9\% versus no action \\
Application & React and WebGL frontend, FastAPI backend \\
Prediction modes & Single order, CSV batch, and public API \\
Deployment & Docker Space on Hugging Face \\
Automation & Two workflows: app deployment and MLOps pipeline \\
Reproducibility & Identical results across three independent runs \\
\bottomrule
\end{tabular}
\end{table}

\section{Reflection}

Four lessons stand out.

\textbf{The most valuable feature did not exist in the data.} Distance had to be
constructed by joining a million GPS readings to both ends of every order.
Feature engineering, not model selection, produced the largest single gain.

\textbf{A model can rank well and still be wrong.} Discrimination and calibration
are independent properties. Our model achieved good ranking while printing
probabilities inflated fourfold, and only measuring calibration exposed it.

\textbf{Deployment is where the real defects appear.} The model trained without
incident. The genuine problems were binary file handling, LFS pointer
resolution, metadata validation, four incomplete dependency declarations,
platform rate limiting, and credential hygiene. Not one of them was a modelling
problem.

\textbf{A working development machine proves almost nothing.} Four separate
faults were invisible locally because a package happened to be installed as an
indirect dependency. Only a clean install exposes an incomplete requirements
file, and a continuous-integration job is exactly a clean install performed
automatically on every change.

\section{Future Work}

\textbf{Seller reliability history} would likely be the single most valuable
addition, computed with time-aware aggregation so that no order sees information
from its own future.

\textbf{Scheduled retraining} would keep the model aligned with changing courier
performance, using the existing workflow extended with a periodic trigger.

\textbf{Monitoring} for input drift and prediction-distribution shift would
detect degradation before it reaches customers.

\textbf{Threshold tuning per region} could account for the wide variation in
baseline reliability between states.

%=============================================================================
\appendix
\chapter{Project Links and Reproduction}
%=============================================================================

\section{Deliverable Links}

\begin{table}[H]
\centering
\caption{Project deliverables}
\begin{tabular}{@{}ll@{}}
\toprule
\textbf{Resource} & \textbf{URL} \\
\midrule
GitHub repository &
\url{https://github.com/moikukreja/olist-delivery-risk} \\
Hugging Face Space &
\url{https://huggingface.co/spaces/samuelalex37/olist-delivery-risk} \\
Registered dataset &
\url{https://huggingface.co/datasets/samuelalex37/olist-delivery-risk-data} \\
Registered model &
\url{https://huggingface.co/samuelalex37/olist-delivery-risk-model} \\
API documentation &
\url{https://samuelalex37-olist-delivery-risk.hf.space/docs} \\
Original dataset &
\url{https://www.kaggle.com/datasets/olistbr/brazilian-ecommerce} \\
\bottomrule
\end{tabular}
\end{table}

The Space runs on free hardware and sleeps after forty-eight hours of
inactivity. If it displays ``Starting'', the container is waking rather than
rebuilding and will be available within approximately thirty seconds.

\section{Reproducing the Results}

The raw data is not committed. After downloading the nine CSV files from Kaggle
into \texttt{data/raw/}, the complete pipeline can be reproduced in order:

\begin{lstlisting}[caption={Reproduction sequence}]
python -m venv venv
venv\Scripts\python.exe -m pip install -r requirements-dev.txt

venv\Scripts\python.exe src\01_data_audit.py
venv\Scripts\python.exe src\02_build_dataset.py
venv\Scripts\python.exe src\03_eda.py
venv\Scripts\python.exe src\04_train_models.py
venv\Scripts\python.exe src\05_smoke_test.py
venv\Scripts\python.exe src\06_build_app_assets.py

venv\Scripts\python.exe -m uvicorn app:app --port 7860
cd frontend && npm install && npm run build
\end{lstlisting}

A fresh clone must run \texttt{git lfs pull} before use, otherwise the model and
data files will be present only as LFS pointers.

\section{Running the MLOps Pipeline}

The four stages can be run individually from a machine holding a Hugging Face
token with write access to the dataset, model and Space repositories:

\begin{lstlisting}[caption={Running the pipeline stages by hand}]
pip install -r mlops/requirements.txt

venv\Scripts\python.exe mlops\data_register.py
venv\Scripts\python.exe mlops\prep.py
venv\Scripts\python.exe mlops\train.py
venv\Scripts\python.exe mlops\hosting.py
\end{lstlisting}

More usually it is run from the repository's Actions tab, under
\textbf{MLOps pipeline}, using \textbf{Run workflow}. Leaving the
\textbf{Promote to Space} option unticked trains and registers a candidate model
without altering the live application, which is the safe default.

The token must be scoped to all three repositories. A token scoped only to the
Space, as used by the application deployment workflow, will authenticate
successfully and then be refused when the pipeline attempts to write the dataset
or the model. The scripts detect this and print the required permission change
rather than a stack trace.

\section{Environment}

\begin{table}[H]
\centering
\caption{Software versions used}
\begin{tabular}{@{}ll@{}}
\toprule
\textbf{Component} & \textbf{Version} \\
\midrule
Python & 3.11.9 \\
scikit-learn & 1.9.0 \\
pandas & 3.0.5 \\
NumPy & 2.4.6 \\
FastAPI & 0.139.2 \\
React & 19.1.0 \\
deck.gl & 9.3.7 \\
Docker base image & python:3.11-slim \\
\bottomrule
\end{tabular}
\end{table}

\projfig{fig02_project_structure}{0.85}{The project directory as organised
during development.}

\projfig{fig28_github_new_repo}{0.95}{Repository creation on GitHub.}

\projfig{fig33_actions_validate_passed}{0.98}{An early workflow run in which
validation passed while deployment failed for want of the secret, demonstrating
the gate operating as intended.}

\end{document}
